% GLSP based multi profile VS Code extension section
% File intended to be included in the main thesis document.
\section{GLSP-based multi-profile VS Code extensions}
\label{sec:glsp-multi-profile}

This section shows how multiple VS Code extension variants are derived from one GLSP-based client codebase, as introduced in Section~\ref{subsec:feature_division}.
The goal is to produce profile-specific products without repository duplication or branch-heavy maintenance.

The section starts with \texttt{profiles.json}, since that manifest is the source of truth for profile selection.
It then follows the manifest into the GLSP client, where modules are assembled and command exposure is restricted by the active build profile.

The main artifacts are \texttt{profiles.json}, \texttt{package.json}, \texttt{build-profile.js}, and the extension webpack configuration.

The flow is straightforward: the manifest defines the server artifact, container modules, and command contribution identifiers; webpack resolves the selected profile, injects compile-time constants, and copies the matching server JAR; activation stores the selected profile in VS Code context and starts the bundled server; command contributions read the same context and expose only compatible interactions.

\subsection{Build-time profiles and manifest structure}

Profiles are declared in a manifest used by the packaging scripts.
Each entry pairs an extension identity with a server JAR, container modules, and command contribution modules.
The split is intentional: container modules contribute GLSP runtime wiring and model-specific services, while command contribution modules provide user-facing commands and editor interactions.
Keeping them separate lets the build reuse shared runtime plumbing while still composing profile-specific behavior.

\lstinputlisting[
  float,
  floatplacement=ht,
  caption={ [\texttt{profiles.json}] Excerpt from profiles.json.},
  label={lst:profiles}
]{code/profiles/profiles.json}

This manifest-driven approach centralizes metadata, making it easier to manage and review.
It also decouples profile definitions from build logic, which improves maintainability and reduces the risk of divergence between products.


\subsection{Build-time injection and packaging}

A VS Code extension is defined by its \texttt{package.json} manifest, which declares contribution points such as commands, custom editors, keybindings, and menus (TODO: https://code.visualstudio.com/api/references/extension-manifest).
Rather than duplicating that manifest across variants, the implemented approach uses one extension with conditional contributions driven by the selected profile.
The webpack configuration resolves the selected profile at build time, injects compile-time constants, and copies the matching server JAR into the extension \texttt{dist} directory.
Each VSIX therefore contains exactly one server artifact.

\lstinputlisting[
  float,
  floatplacement=ht,
  caption={ [\texttt{webpack DefinePlugin}] Define plugin injection (abridged).},
  label={lst:wpdef}
]{code/profiles/webpack_define.js}

Because these values are embedded at build time, runtime branching remains minimal and predictable.

\subsection{Client extension points and \texttt{package.json}}

Profile-specific behavior is gated through context keys such as \texttt{uvl.buildProfile}, evaluated in \texttt{enablement} and \texttt{when} expressions.
This adds a second restriction on top of \texttt{commandContributionIds}: a command must be selected by the manifest and still satisfy the active build profile at runtime.

\lstinputlisting[
  float,
  floatplacement=ht,
  caption={ [\texttt{package.json command}] Command contribution with profile gating.},
  label={lst:pkgcmd}
]{code/profiles/package_command.json}

This approach keeps one manifest while selectively exposing UI elements per profile.
It reduces duplication and keeps variability explicit and reviewable.
In practice, \texttt{buildProfile} is the final guardrail, preventing a contribution from becoming visible when it belongs to another profile.



\subsection{Modular command contributions}

Command registration is organized into small contribution modules.
The build injects an array named \texttt{\_\_UVL\_COMMAND\_CONTRIBUTION\_IDS\_\_}, which the runtime iterates to register only the required contributions.

\lstinputlisting[
  float,
  floatplacement=ht,
  caption={ [\texttt{Contribution registry}] Contribution registry (abridged).},
  label={lst:registry}
]{code/profiles/registry.ts}

This registry pattern reuses common wiring while still allowing profile-specific behavior.
Operationally, it supports RQ3.1 because available user commands are constrained to profile-valid runtime capabilities.

The split between command contributions and container modules also keeps the codebase easier to evolve.
If the GLSP runtime or server-side model wiring changes, the container modules change without forcing command code to move.
If only the user-facing workflow changes, the command modules can be updated independently.

\subsection{Critical appraisal of code references and structure}

The current structure is suitable for producing multiple VSIX artifacts from one codebase.
Centralizing metadata in a manifest improves product governance and lowers divergence risk.
Embedding profile identifiers at compile time reduces runtime complexity and simplifies test design.
Three improvements would strengthen reliability and reproducibility.

\begin{itemize}
  \item Extract the profile selection logic into a testable module with unit tests that validate profile lookup and error handling. Such tests reduce the risk of packaging errors when server artifacts are relocated.
  \item Consider documenting explicit line ranges in the listings or providing a small appendix that reproduces the exact build and packaging commands for reproducibility. This aids verification by external reviewers.
  \item Add an integration test that performs a full build for each profile and asserts that the expected server JAR is present in the packaged \texttt{dist} directory and that the emitted \texttt{\_\_UVL\_BUILD\_PROFILE\_ID\_\_} matches the profile manifest.
\end{itemize}

These recommendations improve packaging reliability and support reproducibility requirements for thesis evaluation.
The repository files substantiating this build flow are \texttt{client/vscode/profiles.json}, \texttt{client/vscode/extension/package.json}, \texttt{client/vscode/build-profile.js}, and \texttt{client/vscode/extension/webpack.config.js}.
